\subsection{逆极限的哈尔测度*}

设 G 是局部紧群（因而完备）。设 \((\mathrm{K}_{\alpha})_{\alpha\in\mathrm{A}}\) 是 G 的紧正规子群的递减有向族，其交为 \(\{e\}\)（所以由 \(K_{\alpha}\) 构成的滤子基收敛到 e）。置 \(G_{\alpha}=G/K_{\alpha}\)；令 \(\varphi_{\alpha}:G\to G_{\alpha}\) 和 \(\varphi_{\beta\alpha}:G_{\alpha}\to G_{\beta}\)（\(\alpha\geqslant\beta\)）为典范同态。于是，逆系统 \((\mathrm{G}_{\alpha},\varphi_{\beta\alpha})\) 的逆极限可与 G 识别，且该逆极限到 \(G_{\alpha}\) 的典范映射可与 \(\varphi_{\alpha}\) 识别（GT, III, §7, No.~3, Prop. 2）。映射 \(\varphi_{\alpha}\) 和 \(\varphi_{\beta\alpha}\) 是适当映射（同上，§4, No.~1, Prop. 1 的 Cor. 2）。本小节始终固定这些假设。

引理 2. --- a) 设 \(f \in \mathcal{K}_{+}(\mathrm{G})\)，S 是包含 Supp \(f\) 的 G 的紧子集，U 是 G 中 S 的开邻域，且 \(\varepsilon > 0\)。存在 \(\alpha \in \mathrm{A}\) 以及函数 \(g \in \mathcal{K}_{+}(\mathrm{G})\)，使 g 在 U 外为零且在 \(\mathrm{K}_{\alpha}\) 的陪集上为常值，并满足 \(|f - g| \leqslant \varepsilon\)。

\begin{enumerate}
\def\labelenumi{\alph{enumi})}
\setcounter{enumi}{1}
\tightlist
\item
  设 \(\mu\) 和 \(\mu'\) 是 \(\mathbf{G}\) 上的两个测度，满足 \(\varphi_{\alpha}(\mu) = \varphi_{\alpha}(\mu')\) 对所有 \(\alpha \in \mathbf{A}\) 成立。则 \(\mu = \mu'\)。
\end{enumerate}

存在 \(\alpha_{1} \in \mathbf{A}\)，使得 \(\mathrm{K}_{\alpha_1}\mathrm{S} \cap \mathrm{K}_{\alpha_1}(\mathrm{G} - \mathrm{U}) = \varnothing\)（GT, II, §4, No.~3, Prop. 4）。扩大 S 并缩小 U 后，可以因此假设 S 和 U 是 \(\mathrm{K}_{\alpha_1}\) 的陪集的并。考虑 S 上满足以下性质的连续数值函数 \(h\)：存在 \(\alpha \geqslant \alpha_{1}\)，使得 \(h\) 在 \(\mathrm{K}_{\alpha}\) 的陪集上为常值。这些函数构成 \(\mathcal{K}(\mathrm{S})\) 的子代数（因为 \((\mathrm{K}_{\alpha})\) 是递减的有向族），包含常数并分离 S 中的点：取 S 中两个不同的点 \(x, y\)；由于 \(\mathrm{K}_{\alpha}\) 的交为 \(\{e\}\)，存在 \(\alpha \geqslant \alpha_{1}\)，使得 \(\varphi_{\alpha}(x) \neq \varphi_{\alpha}(y)\)，于是存在连续数值函数 \(u\)，定义在 \(\varphi_{\alpha}(\mathrm{S})\) 上，使得 \(u(\varphi_{\alpha}(x)) \neq u(\varphi_{\alpha}(y))\)。由 Stone-Weierstrass 定理，存在 \(\alpha \geqslant \alpha_{1}\) 以及在 S 上连续的函数 \(h \geqslant 0\)，使 h 在 \(\mathrm{K}_{\alpha}\) 的陪集上为常值，且 \(|f - h| \leqslant \frac{\varepsilon}{2}\) 在 S 上成立。对于每个 \(t \in \mathbf{R}\)，置 \(\delta(t) = \left(t - \frac{\varepsilon}{2}\right)^+\)，并置 \(h' = \delta \circ h\)。则 \(h'\) 是满足 \(\geqslant 0\) 的函数，在 S 上连续，在 \(\mathrm{K}_{\alpha}\) 的陪集上为常值，且 \(|h - h'| \leqslant \frac{\varepsilon}{2}\) 在 S 上成立，从而 \(|f - h'| \leqslant \varepsilon\) 在 S 上成立。另一方面，\(h'(x) = 0\) 对属于 G 中 S 的边界的 \(x\) 成立，因为此时 \(h(x) \leqslant \frac{\varepsilon}{2}\)。若将 \(h'\) 在 S 的补集上延拓为 0，则得到满足要求的函数 \(g\)，这证明了 a)。

现在设 \(\mu, \mu'\) 是 \(G\) 上的两个测度，满足 \(\varphi_{\alpha}(\mu) = \varphi_{\alpha}(\mu')\) 对所有 \(\alpha \in A\) 成立。令 \(v \in \mathcal{K}(G)\) 为在某个 \(K_{\alpha}\) 的陪集上为常值的函数（其中 \(\alpha \in A\)），于是可写成 \(v = w \circ \varphi_{\alpha}\)，其中 \(w \in \mathcal{K}(G_{\alpha})\)；则 \(\mu(v) = (\varphi_{\alpha}(\mu))(w) = (\varphi_{\alpha}(\mu'))(w) = \mu'(v)\)；由 a) 可知 \(\mu = \mu'\)。

命题 6. --- 对每个 \(\alpha \in \mathbf{A}\)，令 \(\mu_{\alpha}\) 为 \(\mathbf{G}_{\alpha}\) 上的正测度。假设 \(\varphi_{\beta \alpha}(\mu_{\alpha}) = \mu_{\beta}\) 对 \(\alpha \geqslant \beta\) 成立。则存在唯一的正测度 \(\mu\)，定义在 \(\mathbf{G}\) 上，使得 \(\varphi_{\alpha}(\mu) = \mu_{\alpha}\) 对所有 \(\alpha \in \mathbf{A}\) 成立。

唯一性直接由引理 2 b) 得出。现在证明 \(\mu\) 的存在。令 V 为由属于 \(\mathcal{K}(G)\) 且在某个 \(K_{\alpha}\) 的陪集上为常值的函数组成的向量空间（\(\alpha\) 可依函数而定）。由引理 2 a)，V 满足 Ch. III, §1, No.~7, Prop. 9 的条件 (P)：事实上，取 G 中的紧集 K，并选取 \(f \in \mathcal{K}_{+}(G)\)，使 \(f(x) > 0\) 对所有 \(x \in K\) 成立；令 \(a > 0\) 为 \(f\) 在 K 上的最小值；由引理 2 a)，存在函数 \(g \in V \cap \mathcal{K}_{+}(G)\)，使 \(|f - g| \leqslant a/2\)，因此 \(g(x) > 0\) 对所有 \(x \in K\) 成立，条件 (P) 得到验证。取 \(f \in V\)。存在 \(\alpha \in A\)，使 \(f\) 在 \(K_{\alpha}\) 的陪集上为常值。通过取商，\(f\) 定义了函数 \(f_{\alpha} \in \mathcal{K}(G_{\alpha})\)。数 \(\mu(f) = \mu_{\alpha}(f_{\alpha})\) 与 \(\alpha\) 的选择无关：取任意指标 \(\beta\)，使 \(f\) 在 \(K_{\beta}\) 的陪集上为常值；取 \(\gamma \in A\)，使 \(\gamma \geqslant \alpha\)、\(\gamma \geqslant \beta\)；于是 \(f\) 定义函数 \(f_{\beta} \in \mathcal{K}(G_{\beta})\)、\(f_{\gamma} \in \mathcal{K}(G_{\gamma})\)，满足 \(f = f_{\beta} \circ \varphi_{\beta} = f_{\gamma} \circ \varphi_{\gamma}\)；接着 \(f_{\alpha} \circ \varphi_{\alpha\gamma} = f_{\gamma}\)，因此 \(\mu_{\gamma}(f_{\gamma}) = (\varphi_{\alpha\gamma}(\mu_{\gamma}))(f_{\alpha}) = \mu_{\alpha}(f_{\alpha})\)，同理 \(\mu_{\gamma}(f_{\gamma}) = \mu_{\beta}(f_{\beta})\)，从而断言成立。上述成立后，显然 \(\mu\) 是 V 上的线性形式，且 \(\mu(f) \geqslant 0\) 对 \(f \geqslant 0\) 成立。根据 Ch. III, §1, No.~7, Prop. 9，\(\mu\) 可延拓为 G 上的正测度，仍记为 \(\mu\)。有 \(\varphi_{\alpha}(\mu) = \mu_{\alpha}\) 对所有 \(\alpha \in A\) 成立，这正是由 \(\mu\) 的构造得到的。

定义 5. --- 称测度 \(\mu\) 为 \(\mu_{\alpha}\) 的逆极限（或投影极限）。

命题 7. --- 保持命题 6 的记号。如果每个 \(\mu_{\alpha}\) 都是 \(G_{\alpha}\) 上的左（分别为右）哈尔测度，则 \(\mu\) 是 \(G\) 上的左（分别为右）哈尔测度。

例如，设 \(\mu_{\alpha}\) 是左哈尔测度。令 \(s \in \mathbf{G}\)。对每个 \(x \in \mathbf{G}\)，

\[
\left(\varphi_ {\alpha} \circ \boldsymbol {\gamma} (s)\right) (x) = \varphi_ {\alpha} (s x) = \varphi_ {\alpha} (s) \varphi_ {\alpha} (x) = \left(\boldsymbol {\gamma} (\varphi_ {\alpha} (s)) \circ \varphi_ {\alpha}\right) (x);
\]

因此 \(\varphi_{\alpha}(\pmb{\gamma}(s)\mu) = \pmb{\gamma}(\varphi_{\alpha}(s))\mu_{\alpha} = \mu_{\alpha}\)。由引理 2 b)，有 \(\pmb{\gamma}(s)\mu = \mu\)，故 \(\mu\) 是左哈尔测度。

以下设 \(K_{\alpha}\) 不仅是紧的，而且是 G 中的开集。于是 \(G_{\alpha}\) 是离散的；对于 \(\beta \geqslant \alpha\)，\(K_{\alpha}/K_{\beta}\) 是紧的离散群，因而是有限的。群 G 是幺模的（命题 3）。

命题 8. --- a) 设 \(\mu\) 和 \(\mu'\) 是 G 上的两个正测度，且对于每个 \(\alpha\) 以及每个 \(\mathrm{K}_{\alpha}\) 的陪集 C，都有 \(\mu(\mathrm{C}) = \mu'(\mathrm{C})\)。则 \(\mu = \mu'\)。

\begin{enumerate}
\def\labelenumi{\alph{enumi})}
\setcounter{enumi}{1}
\tightlist
\item
  固定一个 \(\alpha_0 \in \mathbf{A}\)。对于每个 \(\alpha \geqslant \alpha_0\)，令 \(n_{\alpha}\) 为有限群 \(\mathrm{K}_{\alpha_0} / \mathrm{K}_{\alpha}\) 的元素个数。存在唯一一个正测度 \(\mu\) 定义在 \(\mathbf{G}\) 上，使得对于每个 \(\alpha \geqslant \alpha_0\)，每个 \(\mathrm{K}_{\alpha}\) 的陪集的测度为 \(n_{\alpha}^{-1}\)。此外，\(\mu\) 是 \(\mathbf{G}\) 上的哈尔测度，并且 \(\mu(\mathrm{K}_{\alpha_0}) = 1\)。
\end{enumerate}

设 \(\mu\) 和 \(\mu'\) 是满足条件 a) 的 G 上两个正测度。于是离散群 \(G_{\alpha}\) 的各点在 \(\varphi_{\alpha}(\mu)\) 和 \(\varphi_{\alpha}(\mu')\) 下具有相同的测度，从而 \(\varphi_{\alpha}(\mu) = \varphi_{\alpha}(\mu')\)，且这对所有 \(\alpha\) 都成立。因此 \(\mu = \mu'\)（引理 2 b)）。

证明 b)。对于每个 \(\alpha \geqslant \alpha_0\)，令 \(\mu_{\alpha}\) 为离散群 \(G_{\alpha}\) 的哈尔测度，使每个点的测度为 \(n_{\alpha}^{-1}\)。设 \(\alpha, \beta\) 满足 \(\alpha \geqslant \beta \geqslant \alpha_0\)。则 \(K_{\beta}/K_{\alpha}\) 有 \(n_{\alpha}/n_{\beta}\) 个元素。因此，\(\varphi_{\beta\alpha}(\mu_{\alpha})\) 是 \(G_{\beta}\) 上每个点的测度为 \(n_{\alpha}^{-1} \cdot \frac{n_{\alpha}}{n_{\beta}} = n_{\beta}^{-1}\) 的测度；换言之，\(\varphi_{\beta\alpha}(\mu_{\alpha}) = \mu_{\beta}\)。于是只需应用命题 6 和 7。

例。--- 设 \(Q_{p}\) 为 p 进数域，即关于 p 进绝对值 \(|x|_{p} = p^{-v_{p}(x)}\) 的 Q 的完备化（GT, IX, §3, No.~2, 例 3）。\(Q_{p}\) 的元素称为 p 进数。仍以 \(|x|_{p}\) 表示 p 进绝对值在 \(Q_{p}\) 上的连续延拓。有

\[
| x + y | _ {p} \leqslant \sup (| x | _ {p}, | y | _ {p})
\]

对于 \(x, y\) 属于 \(\mathbf{Q}\)（同上），因而对于 \(x, y\) 属于 \(\mathbf{Q}_p\) 也成立；此外，如果 \(|y|_p < |x|_p\)，则 \(|x + y|_p = |x|_p\)，因为 \(|x|_p = |(x + y) - y|_p \leqslant \sup(|x + y|_p, |y|_p)\)。若 \((x_n)\) 是 \(\mathbf{Q}_p\) 中趋于 \(x \in \mathbf{Q}_p^*\) 的点列，则 \(|x - x_n|_p < |x|_p\) 且 \(|x - x_n|_p < |x_n|_p\)（当 \(n\) 足够大时），所以 \(|x|_p = |x_n|_p\)。这证明了，对于每个 \(x \in \mathbf{Q}_p^*\)，\(|x|_p\) 是 \(p\) 的幂。

令 \(\mathbf{Z}_p\) 为 \(\mathbf{Z}\) 在 \(\mathbf{Q}_p\) 中的闭包；这是 \(\mathbf{Q}_p\) 的子环；其元素称为 \(p\) 进整数。有 \(|x|_p \leqslant 1\) 对每个 \(x \in \mathbf{Z}_p\) 成立。反之，设 \(x\) 是 \(\mathbf{Q}_p\) 中满足 \(|x|_p \leqslant 1\) 的元素，我们证明 \(x \in \mathbf{Z}_p\)；存在一个序列 \((x_n)\)，其元素属于 \(\mathbf{Q}\) 且趋于 \(x\)，并且 \(|x_n|_p \leqslant 1\) 对足够大的 \(n\) 成立；只需证明 \(x_n\) 属于 \(\mathbf{Z}_p\)（对足够大的 \(n\)）；换言之，我们归结为 \(x \in \mathbf{Q}\) 的情形；此时 \(x = a/b\)，其中 \(b\) 与 \(p\) 互素；对于每个整数 \(n > 0\)，存在 \(b_n' \in \mathbf{Z}\) 和 \(h_n \in \mathbf{Z}\)，使得 \(bb_n' + h_np^n = 1\)，从而 \(x = \frac{abb_n' + ah_np^n}{b} = ab_n' + \frac{ah_np^n}{b}\)，且 \(|x - ab_n'|_p \leqslant p^{-n}\)，于是 \(ab_n'\) 趋于 \(x\)。

由此可知，中心为 0、半径为 \(p^{-n}\) 的闭球，恰好等于中心为 0、半径为 \(p^{-n+1}\) 的开球，即为 \(p^n\mathbf{Z}_p\)。因此拓扑空间 \(\mathbf{Q}_p\) 是零维的，从而是完全不连通的（GT, IX, §6, No.~4）。

下面证明整数 \(0,1,\ldots,p^{n}-1\) 构成 \(\mathbf{Z}_{p}\) 模 \(p^{n}\mathbf{Z}_{p}\) 的一组代表元。首先，有 \(|k-k^{\prime}|_{p}>p^{-n}\)（对于两个这样的整数 \(k\) 和 \(k^{\prime}\)），因此这些整数模 \(p^{n}\mathbf{Z}_{p}\) 的类彼此不同。另一方面，设 \(x\in\mathbf{Z}_{p}\)；存在 \(k\in\mathbf{Z}\) 使 \(|x-k|_{p}\leqslant p^{-n}\)；将 \(p^{n}\) 的适当倍数加到 \(k\) 上，可以设 \(k\in[0,p^{n}-1]\)，且 \(x\) 同余于 \(k\)，模 \(p^{n}\mathbf{Z}_{p}\)。断言得证。这表明 \(\mathbf{Z}_{p}/p^{n}\mathbf{Z}_{p}\) 典范同构于 \(\mathbf{Z}/p^{n}\mathbf{Z}\)。此外可见，\(\mathbf{Z}_{p}\) 是相对紧的，因其完备而紧。由于 \(\mathbf{Z}_{p}\) 是 \(\mathbf{Q}_{p}\) 的开子群，\(\mathbf{Q}_{p}\) 是局部紧的。\(\mathbf{Q}_{p}\) 的拓扑有可数基（GT, IX, §2, No.~9, 命题 16 的推论）。加法群 \(\mathbf{Q}_{p}\) 可与离散群 \(\mathbf{Q}_{p}/p^{n}\mathbf{Z}_{p}\) 的逆极限相识别。

存在唯一一个哈尔测度 \(\alpha\) 定义在加法群 \(\mathbf{Q}_p\) 上，满足 \(\alpha(\mathbf{Z}_p) = 1\)；它称为 \(\mathbf{Q}_p\) 上的归一化哈尔测度。由于 \(\mathbf{Z}_p\) 是 \(p^n\) 个 \(p^n\mathbf{Z}_p\) 的不交陪集的并（\(n\) 为整数且 \(\geqslant 0\)），有 \(\alpha(p^n\mathbf{Z}_p) = p^{-n}\)；类似地，\(\alpha(p^{-n}\mathbf{Z}_p) = p^n\)，所以最终有 \(\alpha(p^n\mathbf{Z}_p) = p^{-n}\) 对每个 \(n \in \mathbf{Z}\) 成立。由命题 8 b)，\(\alpha\) 是 \(\mathbf{Q}_p\) 上唯一的正测度，使得 \(p^n\mathbf{Z}_p\) 的每个陪集（\(n\) 为整数且 \(\geqslant 0\)）的测度为 \(p^{-n}\)。

\(\alpha\) 在 \(\mathbf{Z}_p\) 上的限制显然是 \(\mathbf{Z}_p\) 上的哈尔测度。

